本学では、2025年4月にオープンサイエンス研究開発部門が附属図書館の下部組織として発足した\cite{kgos1}。当部門は、学内外の組織との協力・協調のもとに地域課題解決を目指す「KAGOSHIMAモデル」\cite{kglib1}を推進するための中心的な部局である。
地方大学である本学では、地域社会への貢献が重要な任務であると同時に、その成果をどのように発信あるいは見える化するかもまた、重要な課題となっている。このような社会活動・社会貢献にかかる指標開発や分析は、活動そのものを把握し、より効果的に推進するためにも重要なものとなる。
本部門ではこのための試験的な活動の一環として、ドキュメントから固有表現を抽出しその関係性を整理する仕組みの開発を構想しており（図\ref{fig:kg}）、その第一段階としてAIを導入した地域（地名）を対象とした固有表現抽出（NER）の試みとその評価を行った。
